\chapter{Architekturentscheidungen}\label{app:adr}

%% INTERN-272: KA-162
Dieser Anhang sammelt die tragenden Entwurfsentscheidungen der Arbeit als kompakte
\emph{Architecture Decision Records}\footnote{Architecture Decision Record (ADR): eine kurze, nummerierte
Aufzeichnung einer Entwurfsentscheidung mit Kontext, Entscheidung und Konsequenzen --- eingeführt von
Nygard~\cite{nygard2011adr} als kurze Textdatei in einem Format nach Art eines Alexander-Musters und als
Werkzeug-Format beschrieben von Kopp, Armbruster und Zimmermann~\cite{kopp2018madr}.} (ADRs): je
Entscheidung der Kontext, in dem sie fiel, die Entscheidung selbst und ihre wichtigste Konsequenz. Ziel
ist die schnelle Nachvollziehbarkeit --- \emph{warum} der Apparat so gebaut ist, wie ihn die Kapitel
beschreiben; die ausführliche Herleitung und Begründung liegt jeweils im per \emph{Beleg} referenzierten
Abschnitt des Haupttextes und wird hier nicht wiederholt. Wo Entwurf und Implementierung
auseinanderliegen, weist die ADR den Umsetzungsstand an genau einer Stelle als viertes Feld
aus. Der Bestand umfasst zwanzig Entscheidungen: die zwölf Entscheidungen der Konzeption (ADR-1 bis
ADR-12) und acht Entscheidungen der Identität, der Schnittstellen, des Betriebs und des Baus (ADR-13 bis
ADR-20); Tabelle~\ref{tab:adr-overview} gibt die Übersicht.

\begin{footnotesize}
\begin{longtable}{@{}>{\raggedright\arraybackslash}p{1.3cm}>{\raggedright\arraybackslash}p{6.0cm}>{\raggedright\arraybackslash}p{3.3cm}>{\raggedright\arraybackslash}p{3.8cm}@{}}
\caption{Übersicht der zwanzig Architekturentscheidungen: Kurztitel, belegender Abschnitt und
Umsetzungsstand --- \emph{umgesetzt} (implementiert und belegt), \emph{teilweise} (implementiert, mit
ausgewiesenem offenem Gegenstand) oder \emph{Entwurf} (nur entworfen, nicht implementiert); die Begründungen stehen in den Datensätzen, nicht in
der Tafel}\label{tab:adr-overview}\\
\toprule
ADR & Entscheidung (Kurztitel) & Beleg (Abschnitt) & Umsetzungsstand \\
\midrule
\endfirsthead%
\multicolumn{4}{@{}l}{\tablename\ \thetable{} --- Fortsetzung}\\
\toprule
ADR & Entscheidung (Kurztitel) & Beleg (Abschnitt) & Umsetzungsstand \\
\midrule
\endhead%
\midrule
\multicolumn{4}{@{}r}{\footnotesize Fortsetzung auf der nächsten Seite}\\
\endfoot%
\bottomrule
\endlastfoot%
ADR-1 & C-stabile Modulgrenze statt C++-Modul-ABI & Abschn.~\ref{sec:impl-architecture} & umgesetzt \\
ADR-2 & Compile-Time-Permutation über Typlisten & Abschn.~\ref{sec:axes-impl} & teilweise (Aufschaltung der achtzehnten Achse) \\
ADR-3 & Drei-Ebenen-Modell als die eine Architektur & Abschn.~\ref{sec:anatomy-levels} & umgesetzt \\
ADR-4 & Durchreich-Regel und erzwungene Vollbelegung & Abschn.~\ref{sec:axes-impl} & teilweise (Durchreich-Baustein des Knotentyps) \\
ADR-5 & Schichtenmodell E4--E0 mit festgeschriebenen Verträgen & Abschn.~\ref{sec:anatomy-layers} & umgesetzt \\
ADR-6 & Explizit kodierte Nicht-Erhebung statt Phantomwert & Abschn.~\ref{sec:measurement-system} & umgesetzt \\
ADR-7 & Vier-Repository-Architektur mit orthogonalen Träger-Stufen & Abschn.~\ref{sec:repos} & umgesetzt \\
ADR-8 & Familienweise Adapter-Fallbacks & Abschn.~\ref{sec:axes-impl} & umgesetzt \\
ADR-9 & Mess-Gating in drei Schichten und vier Arbeitsmodi & Abschn.~\ref{sec:impl-pipeline} & teilweise (Replay-Ablauf compare/release) \\
ADR-10 & \texttt{std::map}-Vertrag mit Konformitäts-Gate & Abschn.~\ref{sec:measurement-system} & teilweise (volle \texttt{std::map}-Schnittstelle, Container-Hülle) \\
ADR-11 & Entwurfsmuster-Register mit belegten Rollen & Abschn.~\ref{sec:anatomy-patterns} & umgesetzt \\
ADR-12 & Mess-Regime-Doppel Talos und root-Linux & Abschn.~\ref{sec:eval-method} & teilweise (Talos-Regime) \\
ADR-13 & Vier Träger-Stufen und ein Stempel-Vertrag & Abschn.~\ref{sec:stamp-model} & teilweise (Stempel-Export Hybrid, SOTA-Zweig) \\
ADR-14 & Identitäts-Schnitt: nur Organ-Achsen bilden die \texttt{binary\_id} & Abschn.~\ref{sec:axis-triad} & umgesetzt \\
ADR-15 & Mess-Visitor-Naht (Fläche 3) statt Sidecar & Abschn.~\ref{sec:anatomy-levels} & umgesetzt \\
ADR-16 & Lager als Bestand mit Übersetzungszeit-Schlüsseln & Abschn.~\ref{sec:lager} & teilweise (Lager-Vollausbau) \\
ADR-17 & Optimierungsstufe als Anwender-Wahl, Bau-Welt identitätswirksam & Abschn.~\ref{sec:eval-build-scope} & teilweise (O0/O1-Wahl, Flag-Durchreichung) \\
ADR-18 & Ein offizieller Bauweg als GNU-Vorbau über CMake & Anhang~\ref{app:code} & teilweise (Framework-Export) \\
ADR-19 & Vier CI-Zellen ohne Skip & Abschn.~\ref{sec:impl-contracts} & teilweise (Mini-Pipelines je Träger-Stufe) \\
ADR-20 & XML als einzige Steuerung & Abschn.~\ref{sec:impl-pipeline} & teilweise (Erstellungsmodus, headless Steuerpfad, Schema-Drift) \\
\end{longtable}
\end{footnotesize}

\begin{description}

  \item[ADR-1 --- C-stabile Modulgrenze statt C++-Modul-ABI]\hfill\\
  \textbf{Kontext:} Jede Permutation wird als eigenständige Tier-Binary kompiliert und zur Laufzeit
  geladen; C++23-Module strukturieren zwar die Übersetzungseinheiten, garantieren für sich genommen
  aber keine portable ABI-Stabilität~\cite{wg21_modules_abi}.
  \textbf{Entscheidung:} Eine einheitliche C-Modulgrenze exportiert je Permutation genau ein Lebewesen
  --- eine \texttt{IAnatomyBase}-Instanz --- über sieben gattungsunabhängig benannte Pflicht-Symbole:
  die vier Ur-Symbole ABI-Version, Magic, Factory (\texttt{comdare\_create\_anatomy()}) und Destroy,
  die zwei Identitäts-Symbole Konzept und Gattung, wobei das Konzept nicht als zweites Datum gepflegt,
  sondern übersetzungsstatisch aus der Gattung abgeleitet wird, und das Stempel-Zeilen-Symbol. Die
  Namensgleichheit über alle Gattungen ist das Stabilitäts-Argument: Der Lader kann die Gattung
  über ein konzeptunabhängiges Symbol ermitteln, ohne sie vorab zu kennen. Dahinter liegen zwei Engine-Begriffe: die \texttt{CacheEngine} als
  Bibliothek der Achsen-Bausteine und die \texttt{ExecutionEngine} als Ausführungs- und
  Mess-Infrastruktur mit der Wurzel \texttt{IExecutionEngine}, unter der \texttt{IAnatomyBase} steht;
  die API-Fassade ist \texttt{ICacheEngine}; der Name \enquote{SearchEngine} bezeichnet die ABI-Laufzeit-Sicht
  der Gattung \texttt{SearchAlgorithm}, kein drittes Kettenglied; das gleichnamige Auslagerungs-Modul ist
  deklariert, aber nicht angebunden.
  \textbf{Konsequenz:} Die ABI-Stabilität stellt die C-Grenze her, nicht das Modulsystem (ABI-Major~9);
  der Lader prüft jedes Modul beim Laden fail-closed gegen die Symbol-Pflichtmenge und die
  Identitäts-Ableitung --- ein Modul mit widersprüchlicher Identität oder ohne Stempel-Symbol wird
  abgewiesen, statt als falsche Identität in Stempel, Lager und Messreihe zu wandern.
  \textbf{Beleg:} Abschnitte~\ref{sec:contributions},~\ref{sec:anatomy-patterns}
  und~\ref{sec:impl-architecture}; Abbildung~\ref{fig:one-architecture}.

%% INTERN-272: KA-163
  \item[ADR-2 --- Compile-Time-Permutation über Typlisten statt Laufzeit-Adaption]\hfill\\
  \textbf{Kontext:} Achtzehn Organ-Kompositions-Achsen zählen schon im golden-Referenzkatalog
  $2^{17} = 131072$ Binary-Identitäten je System-Permutation --- eine davon, das Persistenz-Ziel, auf ihren
  In-Memory-Vertreter festgelegt, unter der überholten Annahme zweier Auslegungen je Achse ---, die Obergrenze
  je System-Permutation ist das Produkt der je Achse im XML-Experiment gewählten Auslegungen (erst zur
  Ausführung bekannt), und der rohe Varianten-Produktraum über alle Achsen-Kataloge samt Unter-Achsen liegt
  Größenordnungen darüber;
  der Hot-Path darf weder \texttt{virtual}-Aufrufe noch Laufzeit-Switches enthalten. Die Festlegung ist eine
  Budget-Entscheidung (laut Katalogkopf rund 224~GB und 24~h vollständiger Bau statt rund 448~GB und
  48~h); vorgesehen ist die Variation aller achtzehn Achsen, und die Mächtigkeit eines konkreten Plans
  rechnet der Planer je XML im Director-Walk, nie statisch.
  \textbf{Entscheidung:} Die Permutation läuft übersetzungszeitlich über
  Boost.Mp11-Typlisten\footnote{Boost.Mp11: eine C++11-Metaprogrammier-Bibliothek zur
  übersetzungszeitlichen Manipulation von Datenstrukturen, die Typen enthalten; ihr Listentyp ist
  \texttt{mp\_list}~\cite{dimov_mp11}.}: je Achse eine Gesamt-Liste, per Filter über das
  Aktivierungs-Prädikat zur Enabled-Liste, gebunden an den Kompositions-Slot, dazu concept-geprüfte
  Achsen-Deskriptoren mit dem Inventar \texttt{variants} und eine generische Registry-zu-Baum-Reflektion;
  der golden-Referenzkatalog wählt je Slot die ersten Vertreter der Enabled-Liste --- kein
  Runtime-Tag, kein \texttt{std::variant}. \texttt{std::variant}, \texttt{virtual} und
  Laufzeit-Switches sind in Tier- und Hybrid-Binaries grundsätzlich ausgeschlossen; eine
  statische Prüfung erzwingt das, und der variant-basierte Baustein-Cluster
  (\texttt{baustein\_variants}, \texttt{resolve\_baustein}) ist ein vom Bau isolierter Test-Bestand ohne
  Bibliotheks-Konsumenten.
  \textbf{Konsequenz:} Die Achsen-Abstraktion bleibt eine zero-cost abstraction; in die
  Binary-Identität gehen ausschließlich die achtzehn Organ-Achsen ein (ADR-14).
  \textbf{Umsetzungsstand:} Die Aufnahme der achtzehnten Achse in den Referenzkatalog ist nicht umgesetzt;
  sie verlangt zwei Schritte zugleich --- den Katalog-Grad zwei für das Persistenz-Ziel und den
  Rückschreib-Schalter der Persistenz-Achse.
  \textbf{Beleg:} Abschnitte~\ref{sec:contributions},~\ref{sec:anatomy-patterns},~\ref{sec:axes-impl}
  und~\ref{sec:mess-chain}; Abbildung~\ref{fig:patterns}.

  %% INTERN-272: KA-164
  \item[ADR-3 --- Drei-Ebenen-Modell als die \emph{eine} Architektur]\hfill\\
  \textbf{Kontext:} Anatomie-Metapher und technische Engine-Sicht drohen zu parallelen
  Modellbäumen auseinanderzulaufen.
  \textbf{Entscheidung:} Es gibt genau eine Hierarchie ohne Parallel-Bäume --- Konzept (Ebene~1,
  Mess-Surface: die Prüf-Dock-Konzepte Map, Container und Graph sowie das vierte Konzept
  HeuristikAdapter, das kein eigenes Dock hat, weil seine Binaries per \texttt{genus()} die Ziel-Gattung
  melden und am Dock des Ziels andocken) $\to$ Gattung (Ebene~2, das Tier-Interface mit festem
  Achsen-Satz; die sechs Gattungen \texttt{SearchAlgorithm},
  \texttt{Set}, \texttt{Sequence}, \texttt{Adapter}, \texttt{View} und
  \texttt{FunctionInterfaceReroute}) $\to$ Genus (Ebene~3, die Implementierung des
  Gattungs-Interfaces, z.\,B.\ ART, HOT, START) $\to$ Achsen-Organe (Ebene~4, eigene Klassifikation)
  --- unter \texttt{IExecutionEngine} als gemeinsamer messbarer Wurzel; ein eigenes Konzept entstünde erst, wenn sich der
  Basis-Achsen-Satz grundlegend unterschiede.
  \textbf{Konsequenz:} Die Gattung \texttt{SearchAlgorithm} unter dem Konzept Map und ihre
  ABI-Laufzeit-Sicht (\enquote{SearchEngine}) sind zwei Sichten desselben Lebewesens, kein zweites
  Modell.
  \textbf{Beleg:} Abschnitte~\ref{sec:anatomy-model},~\ref{sec:anatomy-levels}
  und~\ref{sec:impl-architecture}; Abbildungen~\ref{fig:three-levels} und~\ref{fig:genera}.

%% INTERN-272: KA-165
  \item[ADR-4 --- Durchreich-Regel und erzwungene Vollbelegung]\hfill\\
  \textbf{Kontext:} Für nicht-baumartige Verfahren (etwa die Hash-Gegenprobe) sind baum- und
  trie-spezifische Achsen wie Pfadkompression oder Knotentyp nicht sinnvoll belegbar.
  \textbf{Entscheidung:} Solche Achsen fallen nicht weg, sondern werden als Durchreich-Varianten
  (\texttt{None}) belegt --- ein Durchreich-Wert ist ein Algorithmus mit leerem Pfad, kein Weglassen;
  implementiert ist der Durchreich-Baustein für die Pfadkompression
  (\texttt{Path\allowbreak{}Compression\allowbreak{}None}), der Knotentyp besitzt keine None-Variante und wird
  regulär mit einem der vier ART-Knotenformate belegt, die ein nicht-baumartiger Baustein nicht nutzt
  (Entwurf; ohne Code-Baustein);
  das Pflicht-Concept \texttt{IsComposition} der Gattung \texttt{SearchAlgorithm} erzwingt
  übersetzungszeitlich die Belegung aller achtzehn Organe (achtzehn Typ-Forderungen und ein
  Organ-Zähler von achtzehn mit Prüfung auf ABI-Bruch), sodass keine Komposition dieser Gattung mit offener
  Achse übersetzt.
  \textbf{Konsequenz:} Nicht-baumartige Verfahren erfordern kein eigenes Konzept, sondern ein
  eingeschränktes Achsen-Subset innerhalb derselben Gattung.
  \textbf{Beleg:} Abschnitte~\ref{sec:sota-instances},~\ref{sec:anatomy-levels}
  und~\ref{sec:axes-impl}.

  %% INTERN-272: KA-166
  \item[ADR-5 --- Schichtenmodell E4--E0 mit festgeschriebenen Verträgen und Contract-Tests]\hfill\\
  \textbf{Kontext:} Bei einem Produktraum, der jede vollständige Materialisierung ausschließt, sind
  Änderungen über mehrere Ebenen zugleich praktisch nicht abnehmbar; deshalb wird je Ebene einzeln
  fertiggestellt.
  \textbf{Entscheidung:} Die Experiment-Pipeline gliedert sich in die fünf Schichten E4--E0 (E4
  Definition und Auswertung, E3 Permutations-Baum, E2 Tier-Binaries, E1 Laufzeit-Konfiguration, E0
  Infrastruktur) und den Querschnitt~M\@; die E-Kennungen sind eine Konzept-Ordnung mit je einem
  Code-Ort. Jede Schicht wird einzeln fertiggestellt, mit explizit festgeschriebenem Interface-Vertrag zur
  Nachbar-Ebene (die Festschreibung des Vertrags der Schicht E1 im ABI-Adapter: Signaturen, Datenlayout
  und Fähigkeiten festgeschrieben, Semantik-Änderungen nur additiv), und gegen einen übersetzungszeitlichen
  Ersatz (Fake) der Nachbar-Ebene contract-getestet.
  \textbf{Konsequenz:} Jede Ebene ist einzeln abnehmbar, bevor die nächste beginnt; die Grenze
  E3$\to$E2 ist ohne Modul-Bau gegen einen Fake der Schicht E2 contract-getestet, geprüft wird allein
  die statische Baum- und Sichtstruktur samt ihrer Mixed-Radix-Bijektion.
  \textbf{Beleg:} Abschnitte~\ref{sec:anatomy-layers} und~\ref{sec:impl-contracts};
  Abbildungen~\ref{fig:e-layers} und~\ref{fig:impl-layers}.

  %% INTERN-272: KA-167
  \item[ADR-6 --- Explizit kodierte Nicht-Erhebung statt Phantomwert]\hfill\\
  \textbf{Kontext:} Nicht wirklich erhebbare Werte (Zähler ohne Hardware-Zugriff, Achsen ohne
  tatsächlichen Laufzeit-Parameter) laden zu plausiblen Phantomwerten ein, die die Heuristik-Ableitung
  systematisch vergiften würden.
  \textbf{Entscheidung:} Das System meldet explizit ungültig oder leer statt eines Phantomwerts und
  unterscheidet je Messwert drei Zustände: gültiger Messwert, explizit kodierte Nicht-Erhebung (nicht
  anwendbar oder Quelle nicht verfügbar) und Messfehler\footnote{Das Code-Vokabular
  \texttt{honest-0}/\texttt{honest-100}/\texttt{honest-empty} ist der Implementierungsname dieses
  Grundsatzes (Anhang~\ref{app:glossary});
  methodisch ist die Nicht-Erhebung eine ausgewiesene Klasse fehlender Werte im Sinne der
  Missing-Data-Literatur~\cite{graham2009missing}, nie ein imputierter Wert.}. Je Ereignis reist ein
  Gültigkeits-Flag mit, sodass eine gemessene Null ein Messergebnis bleibt; die Wall-Clock-Quelle bietet nur
  Mittelwert und Durchsatz an und verweigert es, ihren Mittelwert als Perzentil zu etikettieren, und
  eine negative Dauer ist ein Messfehler, kein Nullwert; ein PMC-Katalog ohne Modell-Eintrag liefert
  \enquote{nicht anwendbar}, nie eine geratene Kodierung; die Generatoren erzeugen für eine Messreihe
  ohne gültige Zeile keine Ausgabe statt einer leeren.
  \textbf{Konsequenz:} Lücken bleiben sichtbar und gezielt schließbar; der Grundsatz ist in den
  Bausteinen einprogrammiert statt nur dokumentiert.
  \textbf{Beleg:} Abschnitte~\ref{sec:measurement-system},~\ref{sec:impl-system-axes}
  und~\ref{sec:axis-triad}.

  %% INTERN-272: KA-168
  \item[ADR-7 --- Vier-Repository-Architektur mit orthogonalen Träger-Stufen]\hfill\\
  \textbf{Kontext:} Mess-Apparat, Prüfling, Experiment-Definition samt Auswertung und das Manuskript
  haben getrennte Verantwortlichkeiten und Lebenszyklen; die Messwerte werden versioniert direkt in das
  Manuskript eingebaut, weshalb das Manuskript ein eigenes Repositorium mit eigener Pipeline ist.
  \textbf{Entscheidung:} Vier Repositorien --- \texttt{comdare-cache-engine} (\emph{wie} gemessen
  wird; Bausteine-Bibliothek und Builder im Pitchfork-Quellbaum-Layout, Anhang~\ref{app:glossary}),
  \texttt{comdare-prt-art} (der Prüfling), das Projekt-Repository
  \texttt{probst-Diplomarbeit-cache-engine} (\emph{was} durch den Anwender getestet wird:
  Experiment-Definition in XML und XSD, Mess-Treiber, Auswertungs-Werkzeugkette) und das
  Manuskript-Repository mit eigener CI, in das der Vorwärts-Kanal die Generator-Ausgabe als Tabellen-
  und Diagramm-Dateien per Commit einlegt ---; orthogonal zu dieser Aufteilung liegen die vier
  Träger-Stufen in Bau-Reihenfolge Planer $\to$ CacheEngineBuilder $\to$ Tier-Binary $\to$
  Hybrid-Tier-Binary, wobei Stufe $N+1$ nur an Stufe $N$ hängt (Verarbeitungskette: CEB-Prüf-Dock $\to$
  [optionaler Hybrid-Mux] $\to$ Tier-Binary); Planer und CEB sind zwei Binaries --- der Planer
  \texttt{comdare-experiment-planner} und der Mess-Treiber \texttt{comdare-messung-driver} ---, und
  die zwei Rollen erscheinen als Planer-Zeile und CEB-Vertragszeile im Selbst-Stempel des
  Mess-Treibers (\texttt{--version}) in einem Treiber-Binary; den
  Mess-Aggregator führt der \texttt{ExperimentDriver} host-seitig.
  \textbf{Konsequenz:} Builder und CacheEngine teilen sich das Engine-Repository --- der Builder als
  ausführbare Datei über seiner Builder-Bibliothek, die CacheEngine als Bibliothek ---; die Rollen
  bleiben klar getrennt, und der Prüfling ist keine Stufe, sondern der quer liegende Lieferant von
  Organ-Bausteinen.
  \textbf{Beleg:} Abschnitt~\ref{sec:repos} und Anhang~\ref{app:code};
  Abbildung~\ref{fig:traeger-stufen}.

  \item[ADR-8 --- Familienweise Adapter-Fallbacks für eigenständig lauffähige Builds]\hfill\\
  \textbf{Kontext:} Fremd-Code (Allokator-Vendors, SOTA-Bausteine) ist nicht auf jeder Plattform
  und in jedem Build verfügbar oder aktiviert.
  \textbf{Entscheidung:} Übersetzungszeitliche Flag-Familien schalten Adapter zu oder ab --- je
  Vendor ein Familien-Schalter, dessen Nutzung aus Freigabe und Verfügbarkeit folgt, hinter 23
  Vendor-Shim-Headern ---, jeweils mit sicherem, familienweise gewähltem Fallback: Die
  Allokator-Adapter ersetzen den fehlenden Vendor in \texttt{if constexpr}-Zweigen durch die portable
  ausgerichtete Allokation (\texttt{portable\_aligned\_alloc}), der Vendor-Shim liefert für den nie
  ausgeführten Zweig nur Forward-Stubs, und der Achsen-Wert \texttt{StdMalloc} ist ein eigener
  Baustein, keine Notlösung; Datenstruktur-Adapter fremder Suchverfahren schalten über einen
  Familien-Schalter und besitzen eine eigenständige Re-Implementierung als Körper, während die Bindung an
  den Original-Code über den Bibliotheks-Bau mit dem Original-Compiler läuft.
  \textbf{Konsequenz:} Der Build bleibt ohne jede Fremd-Aktivierung eigenständig lauffähig.
  \textbf{Beleg:} Abschnitt~\ref{sec:axes-impl}.

  %% INTERN-272: KA-169
  \item[ADR-9 --- Mess-Gating in drei Schichten und vier Arbeitsmodi]\hfill\\
  \textbf{Kontext:} Mess- und Observer-Kopplungen kosten Zeit und dürfen den Produktiv-Pfad
  nicht belasten.
  \textbf{Entscheidung:} Die Mess-Kopplungen des Laufzeit-Adapters werden über die Master-Option
  \texttt{COMDARE\_MEASUREMENT\_MODE} geschaltet, die \texttt{COMDARE\_MEASUREMENT\_ON} definiert und die
  der Release-Modus erzwungen abschaltet; die schaltbare Fläche zerfällt in drei trennbare Schichten
  --- Basis-Zeit, Observer-Zähler und Feinkorn-Segment-Timer je Achse ---, die Abschaltungsstufen der
  Messung. Die Ablaufmethodik kennt vier Arbeitsmodi derselben Permutation, als Registry und nicht als
  Etikett: \texttt{build} (bauen ohne Messen), \texttt{measure} (deterministischer
  Ein-Thread-Messlauf, der einzige Modus mit Messung), \texttt{compare} (liest das Messwertlager, die
  Stufe vor \texttt{release}) und \texttt{release} (Voll-Optimierung ohne Mess-Instrumentierung); der
  Arbeitsmodus ist eine Unter-Achse des Mess-Realms und geht nicht in die Binary-Identität ein, einen
  Debug-Zustand gibt es nicht (\texttt{--debug} ist ein Flag der Planer-Shell; unter \texttt{--debug} fährt
  das Konformitäts-Gate zusätzlich den Genus-Testsatz, ein indirekter Debug-Modus, der mit dem
  Release-Bau zusammenfällt), und den Mess-Aggregator führt der \texttt{ExperimentDriver} host-seitig.
  \textbf{Konsequenz:} Der Produktiv-Pfad bleibt mess-frei und verursacht keinen Mess-Overhead; wählbar sind
  vier Modi, in der Bau-Semantik verschieden ist allein \texttt{measure} --- alle vier Registry-Zeilen
  führen den Release-Bautyp.
  \textbf{Umsetzungsstand:} Der modus-spezifische Replay-Ablauf von \texttt{compare} und die aus
  Messwerten erzeugte optimale Release-Binary sind vorgesehen, aber nicht implementiert;
  \texttt{compare} ist ein wählbares, validierbares Etikett mit release-gleicher
  Bau-Semantik.
  \textbf{Beleg:} Abschnitte~\ref{sec:repos},~\ref{sec:impl-pipeline} und~\ref{sec:eval-method};
  Abbildungen~\ref{fig:modi-kette} und~\ref{fig:abschaltungsstufen}.

  %% INTERN-272: S-112
  %% INTERN-272: KA-170
  \item[ADR-10 --- \texttt{std::map}-Vertrag als gemeinsamer Nenner mit Konformitäts-Gate]\hfill\\
  \textbf{Kontext:} Strukturell verschiedene Suchalgorithmen sind nur über eine gemeinsame, gut
  verstandene Außen-Schnittstelle fair vergleichbar.
  \textbf{Entscheidung:} Alle Verfahren werden auf den C++-Vertrag der geordneten Map
  \texttt{std::map<Key,\,Value>}~\cite{iso_cpp} abgebildet; die Container-Gattungen werden gegen je ein
  Standard-Orakel gehalten --- \texttt{Set} gegen \texttt{std::set}, \texttt{Sequence} gegen
  \texttt{std::deque}, \texttt{Adapter} gegen \texttt{std::queue}, \texttt{std::stack} und
  \texttt{std::priority\_queue}, \texttt{View} gegen \texttt{std::span} ---, und \texttt{std::vector}
  dient als Inner-Substrat der Gattung \texttt{Adapter}. Ein Konformitäts-Gate prüft diesen Vertrag vor
  jeder Messung in der bindenden Reihenfolge Import $\to$ Gate $\to$ Messen: die siebzehn in
  Anhang~\ref{app:interfaces} markierten \texttt{std::map}-Operationen (zwölf Kern-Prüfblöcke,
  \texttt{std::map} als Orakel), mit Konformitäts-Quoten als Ergebnisform; die fünf Antriebs-Operationen
  der Modulgrenze (\texttt{insert}, \texttt{lookup}, \texttt{erase}, \texttt{clear}, \texttt{size}) sind
  die ABI-Primitive, über denen der host-seitige Adapter diesen Vertrag komponiert, nicht der
  Prüfumfang. Unter \texttt{--debug} fährt das Gate zusätzlich den Standard-Testsatz des jeweiligen
  Tier-Interfaces: den generischen Gattungs-Teil (gültig für alle Genera) und den Genus-speziellen Teil
  (die gegebenenfalls erweiterten Spezial-Interfaces); dieser indirekte Debug-Modus fällt mit dem
  Release-Bau zusammen und wird nicht separat geführt.
  \textbf{Konsequenz:} Nur Vergleichbares wird verglichen; die Achsen beschreiben ausschließlich
  das innere Verhalten hinter der festen Schnittstelle.
  \textbf{Umsetzungsstand:} Der Ausbau des Gates auf die volle \texttt{std::map}-Schnittstelle und die
  Gattungs-Iterator-Bibliothek sind nicht implementiert; die Forderung, die Container-Gattungen über eine
  vollständige \texttt{std::vector}-Hülle zu führen, ist nicht erfüllt --- das implementierte Sequence-Orakel ist
  \texttt{std::deque}.
  \textbf{Beleg:} Abschnitte~\ref{sec:measurement-system} und~\ref{sec:impl-contracts}.

  \item[ADR-11 --- Entwurfsmuster-Register mit belegten Rollen]\hfill\\
  \emph{Command als Bündelungs- und Befehlsmuster.}\hfill\\
  \textbf{Kontext:} Der Apparat ist bewusst eine Komposition benannter Entwurfsmuster; unklare
  Muster-Rollen würden die Architektur verwässern.
  \textbf{Entscheidung:} Verwendet werden, je mit benanntem Code-Ort: \emph{Strategy} mit
  CRTP~\cite{coplien1995crtp} je Organ-Achse (auch der Stempel-Vertrag der Träger ist ein
  CRTP-Vertrag, ADR-13), \emph{Abstract Factory} (Hybrid-Docks, Kompositions-Fabrik,
  Prüflings-Registry), \emph{Director} mit \emph{Builder} (Plan-Emission), \emph{Iterator} (der
  lazy durchlaufene Experimentbaum), \emph{Adapter} (Fremd-Code hinter der ABI-Grenze),
  \emph{Observer} (Mess-Snapshots je Achse), \emph{Visitor} (die Mess-Naht am
  \texttt{dlopen}-Übergang, ADR-15), \emph{Chain of Responsibility} (Hardware-Erhebungskette und
  Auswahl-Filterkette), \emph{Memento} (Copy-on-Write im ABI-Adapter, Zustand je Tier),
  \emph{Command} und \emph{Facade} --- dieselben Rollen wie das Muster-Register in
  Abschnitt~\ref{sec:impl-hybrid-lager} (Abbildung~\ref{fig:muster}), benannt nach dem Katalog von
  Gamma et al.~\cite{gamma1994patterns} und als Tupel von Achsen-Policies
  komponiert~\cite{alexandrescu2001modern}. Die Muster \emph{Composite} und \emph{Decorator} gehören nicht dazu:
  Die Mess-Hüllen der Achsen halten keine Komponenten-Instanz und kein Voll-Interface und sind
  ausdrücklich kein Decorator, und ein Composite im Sinne des Katalogs hat keine Entsprechung im Code.
  \emph{Nicht} verwendet werden im Hot-Path und in den Organ-Achsen \emph{Singleton}, \emph{Flyweight}
  und \emph{Bridge}; host-seitige Registries nutzen Meyers-Singletons\footnote{Meyers-Singleton: die
  eine Instanz als funktionslokales statisches Objekt, das beim ersten Aufruf initialisiert
  wird~\cite{meyers2005effective}.}, die Katalog-Erzeugung Flyweight-Tabellen je Achse, und
  \enquote{Bridge} ist im Code nur eine umgangssprachliche Bezeichnung für Verbindungsstücke
  (ABI-Adapter, CI-Bridge-Jobs, Sandy Bridge), kein Muster. Das \emph{Command}-Muster
  trägt die übersetzungszeitliche Bündelung von Achsen-Rekombinationen (die System-Komplex-Achse als
  Command-Wrapper, der Hub der externen Werkzeuge als Invoker, die Command-Basis der Achsen im
  Querschnitt~M) und den typisierten Steuer-Befehlskanal; der Befehls-Kanal des Hybrid-Modus ist eine
  dieser Rollen.
  \textbf{Konsequenz:} Jedes Muster trägt benannte Rollen mit Code-Ort in der Verkabelung des Apparats
  und wird übersetzungszeitlich statt laufzeit-polymorph ausgeführt.
  \textbf{Beleg:} Abschnitte~\ref{sec:anatomy-patterns},~\ref{sec:heuristics-modes}
  und~\ref{sec:impl-hybrid-lager}.

  %% INTERN-272: KA-171
  \item[ADR-12 --- Mess-Regime-Doppel: immutables Talos und root-Linux]\hfill\\
  \textbf{Kontext:} Produktions-Treue und Mess-Tiefe stehen im Konflikt: Ein produktionsnahes
  System verwehrt gerade den Hardware-Zähler-Zugriff, den tiefe Messungen brauchen.
  \textbf{Entscheidung:} Auf den Produktions-Zielmaschinen wird jede Messung unter zwei
  Betriebssystem-Regimes erhoben: einem immutablen Betriebssystem (Talos\footnote{Talos Linux: das
  unveränderliche, API-gesteuerte Betriebssystem für Kubernetes von Sidero
  Labs~\cite{siderolabs_talos}.}) als Produktiv-Spiegel und einem root-Linux mit vollem
  Hardware-Zähler-Zugriff (\texttt{perf}~\cite{demelo2010perf,linux_perf_event_open} und die
  modellspezifischen Register, MSR).
  \textbf{Konsequenz:} Die zeit- und observer-basierten Kategorien laufen unter beiden Regimes
  identisch und sichern die Vergleichbarkeit; die PMC-Kategorien werden nur unter dem
  privilegierten Regime erhoben.
  \textbf{Umsetzungsstand:} Die Entscheidung gilt; umgesetzt ist die root-Linux-Hälfte: Die berichteten
  Messungen stammen sämtlich vom root-Linux-Regime der beiden Produktionsmaschinen --- prod1 mit der
  AMD-Mikroarchitektur Zen~5 und prod2 mit dem Intel Core i9-12900K (Alder Lake) ---; das
  Talos-Regime ist vorgesehen, aber nicht gemessen. Für Intel enthält der PMC-Rohereignis-Katalog
  keine Belegung, für Zen~5 ist sie durch Messung gegengeprüft.
  \textbf{Beleg:} Abschnitte~\ref{sec:eval-method} und~\ref{sec:contributions}.

%% INTERN-272: KA-172
  \item[ADR-13 --- Vier Träger-Stufen und ein Stempel-Vertrag]\hfill\\
  \emph{Der Stempel als Identität, Lager-Schlüssel, Wiedererkennungs-Kriterium und Einordnung.}\hfill\\
  \textbf{Kontext:} Die Konfiguration wandert durch vier Träger-Stufen (Planer, CacheEngineBuilder,
  Tier-Binary und die optional zwischen CEB-Prüf-Dock und Tier-Binary eingeschobene Hybrid-Tier-Binary);
  jede Stufe braucht eine ABI-stabile, übersetzungszeitliche Identität, in die kein Messwert eingeht.
  \textbf{Entscheidung:} Ein einziger CRTP-Vertrag bindet alle vier Träger: sieben Stempel-Interfaces
  (Version, Mess-Zeile, System-Zeile, Organ-Zeile, Fingerprint, Gesamt-Stempel, Angeschlossene) mal
  vier Träger ergeben 28 dreiwertige Zulassungszellen (Tabelle~\ref{tab:stempel-matrix}); die Absenz
  der Organ-Zeile beim CEB ist die statische Sperre \enquote{der CEB führt keine Organ-Achsen}, und
  die Prüfung läuft als \texttt{static\_assert} unter jeder Erbin. Der Planer führt einen
  SHA-256-Fingerprint aus Version, ISA und Betriebssystem, die Tier-Binary einen SHA-512-Fingerprint
  über elf Glieder (Format~6), darunter die drei Realm-Zeilen, das Werkzeugketten-Glied und den Hash
  der Overlay-Quellmenge (Konkatenation in fester Ordnung, einmal SHA-512); der CEB führt einen
  SHA-512-Provenienz-Fingerprint, der kein Wiederverwendungs-Kriterium ist. Die Klammer-Grammatik der
  Stempel-Einträge ist an einer Stelle definiert (\texttt{c\{p.e\}} ist \texttt{c} mit den Untergliedern \texttt{p} und
  \texttt{e}, nie flach), und die Klammer-Tiefe ist die Meta-Meta-Ebene des Eintrags.
  \textbf{Konsequenz:} Der Stempel erfüllt fünf Rollen --- Identität, Cache-Schlüssel, Lager-Schlüssel,
  Wiedererkennungs-Kriterium und Einordnung; der Selbst-Stempel des Treiber-Binaries (\texttt{--version}) führt nur
  die Einordnung und ist weder Cache-Schlüssel noch Bau-Eingang.
  \textbf{Umsetzungsstand:} Der Registry-Wrapper-Stempel des SOTA-Zweigs und der Stempel-Export der
  Hybrid-Stufe sind nicht implementiert; die Stempel-Identität in jedem Diagramm und jeder Tabelle des
  Anhangs ist ebenfalls nicht implementiert.
  \textbf{Beleg:} Abschnitte~\ref{sec:stamp-model},~\ref{sec:anatomy-levels},~\ref{sec:mess-chain},~\ref{sec:impl-contracts}
  und~\ref{sec:eval-identity}; Abbildung~\ref{fig:traeger-flaechen}.

  %% INTERN-272: KA-173
  \item[ADR-14 --- Identitäts-Schnitt: nur Organ-Achsen bilden die \texttt{binary\_id}]\hfill\\
  \emph{System- und Mess-Achsen sind identitätsneutral.}\hfill\\
  \textbf{Kontext:} Drei Achsen-Realms (Organ, System, Mess) führen je eine eigene Registry mit
  generierter XML\@; gingen Mess- oder System-Werte in die Identität ein, entstünden Doppel-Binaries
  derselben Permutation und ein Lager, das dieselbe Binary mehrfach hält.
  \textbf{Entscheidung:} Die \texttt{binary\_id} ist der Pfad aus achtzehn \texttt{achse=wert}-Segmenten
  in kanonischer Slot-Reihenfolge (die Organ-Zahl achtzehn ist an genau einer Stelle definiert); alle Mess-Achsen
  und alle System-Achsen führen in ihren Registries \texttt{binary\_id = never}; Unter-Achsen sind
  Laufzeit-Permutation und stehen in keiner Bau-Job-Legende; der Arbeitsmodus ist identitätsneutral.
  Die Bau-Welt reist stattdessen als Fingerprint-Glied mit --- die System-Zeile und das
  Werkzeugketten-Glied mit \texttt{vendoropt} (ADR-17).
  \textbf{Konsequenz:} Lager-Schlüssel und Regressions-Detektor (Referenzkatalog $2^{17}$ mit
  CRC-64-Prüfsumme) sind organ-stabil; dieselbe Organ-Permutation gegen zwei Bau-Welten ist ein
  verschiedenes Artefakt mit verschiedenem Fingerprint.
  \textbf{Beleg:} Abschnitte~\ref{sec:axis-triad},~\ref{sec:stamp-model}
  und~\ref{sec:impl-system-axes}; Tabelle~\ref{tab:axes-overview}.

  %% INTERN-272: KA-174
  \item[ADR-15 --- Mess-Visitor-Naht (Fläche 3) statt Sidecar]\hfill\\
  \emph{Konzept- und Gattungs-Interfaces bleiben unverändert.}\hfill\\
  \textbf{Kontext:} Ein Beobachter, der von außen Felder aus einer Tier-Binary abholt, misst die
  Nachbarschaft und nicht den Gegenstand: Was gemessen wird, entscheidet dann der Host durch die Wahl
  der Felder; Mess-Felder in den fachlichen Interfaces würden die Konzept- und Gattungs-Verträge
  verändern.
  \textbf{Entscheidung:} Jede Träger-Stufe bietet drei Interface-Flächen: Fläche~1 das
  Gattungs-Interface (Laufzeit, Abstract Factory), Fläche~2 den Stempel (Übersetzungszeit, ABI-stabil),
  Fläche~3 den Mess-Durchstich; die Messung quert als eigene Naht
  (\texttt{tier\_measure\_accept} mit \texttt{IMessVisitor}), ein Visitor am
  \texttt{dlopen}-Übergang mit doppeltem Dispatch~\cite{gamma1994patterns}; je nachdem, ob die
  Messeinrichtung in CEB und Tier-Binary aktiviert ist, wird am Interface ein Mess-Visitor übergeben
  oder nicht.
  \textbf{Konsequenz:} Die fachlichen Interfaces enthalten keine Mess-Felder; die Gate-Kombination CEB
  aus, Tier an ist durch einen Deckungs-Test belegt; die Mess-Gate-Belegung reist als eigenes
  Fingerprint-Glied mit eigenen Positionen für die Hybrid-Ebene.
  \textbf{Beleg:} Abschnitte~\ref{sec:anatomy-levels},~\ref{sec:impl-pipeline}
  und~\ref{sec:measurement-system}; Abbildung~\ref{fig:traeger-flaechen}.

  %% INTERN-272: KA-175
  \item[ADR-16 --- Lager als Bestand mit Übersetzungszeit-Schlüsseln]\hfill\\
  \emph{Messung immer an, im Lager nur Fehlendes nachmessen.}\hfill\\
  \textbf{Kontext:} Ein Referenzraum von $2^{17}$ Kompositionen je System-Permutation im golden-Referenzkatalog (laut
  Katalogkopf rund 224~GB und 24~h vollständiger Bau; überholte Annahme zweier Auslegungen je Achse, die
  Obergrenze bestimmt der Planer aus dem XML-Experiment) ist nur mit Wiederverwendung gebauter Binaries und
  gemessener Werte beherrschbar; die Messung soll immer eingeschaltet sein und im Lager nur
  nachmessen, was dort nicht zu finden ist.
  \textbf{Entscheidung:} Bestand~1 hält die Binary unter ihrem SHA-512-Fingerprint; Bestand~2 hält
  die Messwerte unter dem Paar aus Fingerprint der vermessenen Binary und Hardware-Identität der
  messenden Maschine, die Mess-Zelle steht daneben und wird nie in den Digest fusioniert; eine
  vorgefundene Binary wird übersprungen, wenn ihr bestandenes Prüf-Testat inventarisiert ist, sonst gebaut
  oder gemessen; ein Versions-Lock über die Overlay-Quellmenge macht unbemerkte Quelländerungen zum
  CI-Fehler.
  \textbf{Konsequenz:} Die Identität reist unverändert vom Lager über die Mess-Zeile
  (\texttt{permutation\_id}, \texttt{fingerprint}) in Mappe und Anhang.
  \textbf{Umsetzungsstand:} Bestand~1 und~2 und das Schlüssel-Schema sind implementiert; Bestand~3 und~4
  (Synthese-Batches und Lösungen), die Planer-Blattfunktion, die CEB-Vollplatzierung, die
  Hybrid-Ablage, der Commit-Skip, das Rebuild-Flag je Stufe und die XML-Sektion \texttt{publish} bilden
  den Lager-Vollausbau und sind nicht implementiert.
  \textbf{Beleg:} Abschnitte~\ref{sec:lager},~\ref{sec:eval-method},~\ref{sec:stamp-model}
  und~\ref{sec:heuristics-modes}; Abbildungen~\ref{fig:lager-bestaende}
  und~\ref{fig:eval-lager-kreislauf}.

  %% INTERN-272: KA-176
  \item[ADR-17 --- Optimierungsstufe als Anwender-Wahl]\hfill\\
  \emph{Haus-weit O2, O3 auf XML-Wunsch; die Bau-Welt ist identitätswirksam.}\hfill\\
  \textbf{Kontext:} Messreihen über verschieden optimierte Bau-Welten sind nicht vergleichbar;
  die maximale Optimierung muss dem Anwender wählbar bleiben.
  \textbf{Entscheidung:} Die Release-Flags gelten haus-weit mit O2 --- einschließlich Vendor-Bau,
  FetchContent und Einbettung --- per erzwungenem Cache-Wert; O3 wird nur auf XML-Eingangswunsch
  mitgebaut; die Bau-Welt-Stufe geht als Feld \texttt{vendoropt} in das Werkzeugketten-Glied des
  Fingerprints ein (Übersetzungsfehler statt stillem Default); ein Gate verwirft eigene
  Release-Optimierungsflags an der Achse vorbei; die Unter-Achse \texttt{opt\_level} bietet O0 bis
  Ofast mit dem Determinismus-Flag nach IEEE~754~\cite{ieee754_2019}, das allein Ofast verneint.
  \textbf{Konsequenz:} Zwei Baue derselben Permutation gegen eine O2- und eine O3-Welt sind
  verschiedene Artefakte mit verschiedenem Fingerprint; ein stiller Skip ist ausgeschlossen.
  \textbf{Umsetzungsstand:} O0 und O1 als Anwender-Wahl in der XML und die Durchreichung weiterer
  Compiler-Flags aus der XML sind vorgesehen, aber nicht implementiert (Default O2, O3
  überschreibend wählbar).
  \textbf{Beleg:} Abschnitte~\ref{sec:eval-build-scope},~\ref{sec:impl-system-axes}
  und~\ref{sec:complex-axis}.

  %% INTERN-272: KA-177
  \item[ADR-18 --- Ein offizieller Bauweg als GNU-Vorbau über CMake]\hfill\\
  \emph{\texttt{configure.sh}, \texttt{make}, \texttt{make check} und \texttt{make install}.}\hfill\\
  \textbf{Kontext:} Das einzige unter Linux offizielle Verfahren ist \texttt{configure.sh},
  \texttt{make} und \texttt{make install} beziehungsweise \texttt{make check}, und diese drei müssen
  im Wurzelordner liegen; eine Konfiguration darf nur eine Quelle haben.
  \textbf{Entscheidung:} \texttt{configure.sh} ist ein POSIX-sh-Vorbau ohne Python, gebaut wird
  weiterhin mit CMake; das \texttt{Makefile} ist statisch eingecheckt und reicht nur an CMake weiter,
  mit den GNU-Standardzielen~\cite{gnu2026standards} \texttt{all}, \texttt{install}, \texttt{check},
  \texttt{installcheck}, \texttt{uninstall}, \texttt{clean}, \texttt{distclean},
  \texttt{mostlyclean} und \texttt{maintainer-clean} (das Test-Ziel heißt \texttt{check}); die
  Konfiguration ist ausschließlich in \texttt{config.status} und im CMake-Cache abgelegt, und das Makefile
  enthält bewusst keine Vorgabewerte --- ein generiertes Makefile wäre eine zweite Konfigurationsquelle ohne
  Nutzen.
  \textbf{Konsequenz:} Es gibt keinen zweiten Bauweg neben CMake und keine zweite
  Konfigurationsquelle; der Vorbau ist mit \texttt{sh -n} prüfbar.
  \textbf{Umsetzungsstand:} \texttt{configure.sh} und \texttt{Makefile} liegen im Repositorium; der
  Framework-Export für Konsumenten (\texttt{install(EXPORT)}, \texttt{find\_package}) und die
  Nachlauf-Prüfungen sind nicht implementiert.
  \textbf{Beleg:} Anhang~\ref{app:code}.

  %% INTERN-272: KA-178
  \item[ADR-19 --- Vier CI-Zellen ohne Skip]\hfill\\
  \emph{gcc und clang je Release und Debug; ein Job schlägt fehl oder besteht.}\hfill\\
  \textbf{Kontext:} Zwei Übersetzer finden verschiedene Fehler; ein stiller Job-Skip oder ein
  tolerierter Fehler täuscht einen bestandenen Lauf vor.
  \textbf{Entscheidung:} Die Test-Suite läuft in vier Zellen --- gcc und clang jeweils als Release-
  und Debug-Bau --- ohne Änderungs-Skip; die Zelle ist die Warnung, der Job schlägt fehl;
  \texttt{allow\_failure} führt nur zwei deklarierte Ausnahmen --- den arm64-Smoke-Bau ohne
  verfügbaren Runner und den manuellen Relock-Wartungsjob ---, keine Bau- oder Test-Zelle; je
  Übersetzer und Bautyp gibt es einen Objekt-Cache-Schlüssel mit genau einem Schreiber; der
  Übersetzer-Mindeststand ist g++~15.3; der golden-Referenzkatalog $2^{17}$ mit CRC-64-Prüfsumme ist
  der Regressions-Detektor mit festgeschriebener Kardinalität.
  \textbf{Konsequenz:} Akzeptiert wird nur ein Lauf ohne fehlgeschlagene Tests in allen vier Zellen ohne
  Ausnahmeliste --- ein Gate im Sinne der Continuous Delivery~\cite{humble2010cd} ---, und eine
  Test-Zahl wird nur mit dem Lauf genannt, der sie erzeugt hat.
  \textbf{Umsetzungsstand:} Die Mini-Pipeline je Träger-Stufe (vier Stufen mal gcc und clang mal Release
  und Debug) ist vorgesehen, aber nicht implementiert.
  \textbf{Beleg:} Abschnitt~\ref{sec:impl-contracts}.

  %% INTERN-272: KA-179
  \item[ADR-20 --- XML als einzige Steuerung]\hfill\\
  \emph{Die XML trägt Was, Wo und Wann, nie das Wie.}\hfill\\
  \textbf{Kontext:} Der Apparat wird aus einem deklarativen Experiment-Profil betrieben; Begriffe
  werden nicht übersetzt, sondern müssen konform sein --- wäre eine Übersetzung nötig, ist das eine
  Regression und ein Übersetzungszeit-Fehler.
  \textbf{Entscheidung:} Zwei Schemata beschreiben Profil und Messreihe
  (\texttt{experiment\_schema.xsd}, \texttt{messreihe\_v32\_schema.xsd}); die Parse-Naht bricht mit Fehlermeldung
  ab, die Validierung ist ein eigener Baustein, und die CI prüft die Wohlgeformtheit; die XML beschreibt
  Was, Wo und Wann, nie das Wie; die Achsen-Registries werden als XML generiert und sind damit
  maschinenlesbare Abbilder; der Planer ist ein Subkommando-Werkzeug (\texttt{validate},
  \texttt{plan}, \texttt{cache-key}, \texttt{fingerprint}, \texttt{status}, \texttt{check-size},
  \texttt{simulate}, \texttt{version}, \texttt{help}) und der Shell-Steuerpfad; als zweiter Steuerpfad ist eine
  API-Kommandozeile mit headless Dienst vorgesehen, der den Planer als Zelle auf Kubernetes oder auf
  bare metal bedient, sowie ein kommandozeilengeführter Erstellungsmodus für die Anwender-XML, in dem
  nicht verfügbare Optionen nicht wählbar sind.
  \textbf{Konsequenz:} Es gibt keinen Nebenweg neben der XML\@; die SOTA-Profile sind als XML
  ausführbar.
  \textbf{Umsetzungsstand:} Der headless Steuerpfad und der Erstellungsmodus sind nicht implementiert;
  die Hybrid-Sektion ist geparst, aber im Schema nicht deklariert, die
  Organ-Meta-Meta-Sektion deklariert, aber nicht gelesen, und die \texttt{publish}-Sektion fehlt.
  \textbf{Beleg:} Abschnitte~\ref{sec:mess-chain} und~\ref{sec:impl-pipeline};
  Abbildung~\ref{fig:pipeline}.

\end{description}
